% Adapted from the Star Rover template by Subidit.
% Source: https://github.com/subidit/rover-resume/tree/main/templates/star%20rover
% License: CC BY 4.0 (see TEMPLATE-LICENSE.md)
\documentclass[11pt]{article}

% Star Rover uses Fira Sans upstream. Use LaTeX's built-in sans-serif family so
% the CV also compiles when the optional FiraSans package is unavailable.
\renewcommand{\familydefault}{\sfdefault}
\usepackage[a4paper,margin=1in]{geometry}
\usepackage{xcolor}
\definecolor{accent}{HTML}{141E61}

\setcounter{secnumdepth}{0}
\pdfgentounicode=1

% List formatting
\usepackage{enumitem}
\setlist[itemize]{nosep,left=0pt .. 1.5em}
\setlist[enumerate]{left=0pt..1.5em}
\setlist[description]{itemsep=0pt}

% Title formatting
\usepackage{xhfill}
\usepackage{soul}
\newcommand{\midrule}[1]{%
  \leavevmode
  \xrfill[.5ex]{1pt}[accent]~\so{#1}~\xrfill[.5ex]{1pt}[accent]}

\usepackage{titlesec}
\titlespacing{\section}{0pt}{*4}{*1}
\titlespacing{\subsection}{0pt}{*3}{*0}
\titlespacing{\subsubsection}{0pt}{*0}{*0}
\titleformat{\section}{\large\bfseries\uppercase}{}{}{\midrule}
\titleformat*{\subsection}{\large\bfseries\scshape}

% Star Rover entry helpers
\newcommand{\aux}[1]{$|$ \space {\normalfont\itshape #1}}
\newcommand{\rside}[1]{\hfill {\normalfont\color{accent} #1}}

% Footer
\usepackage{lastpage}
\usepackage{fancyhdr}
\pagestyle{fancy}
\fancyhf{}
\renewcommand{\headrulewidth}{0pt}
\fancyfoot[C]{\small\color{gray}Pingshi Yu -- Page \thepage\ of \pageref*{LastPage}}

\usepackage{multicol}
\usepackage[bookmarks=false,hidelinks]{hyperref}
\newcommand{\linkicon}{{\color{gray}\footnotesize[link]}}

\begin{document}

% Banner
\begin{center}
  {\fontsize{32}{36}\fontseries{heavy}\selectfont\color{accent}
    PINGSHI (JACOB) YU}\\[5pt]
  London, UK \quad
  \href{mailto:p.yu22@imperial.ac.uk}{p.yu22@imperial.ac.uk} \quad
  \href{https://pingshiyu.github.io}{pingshiyu.github.io}\\[3pt]
  \href{https://github.com/pingshiyu}{GitHub: pingshiyu} \quad
  \href{https://www.linkedin.com/in/jacob-pingshi-yu/}{LinkedIn} \quad
  \href{https://scholar.google.com/citations?user=hbgX7rEAAAAJ}{Google Scholar}
\end{center}

\section{Research Profile}

Final-year Computer Science PhD student studying the theory and practice of reliable software systems: how to find bugs when they occur and how to avoid them in the first place. Research spans compiler testing, programming-language semantics, property-based testing, complexity theory, and automated verification. Supervised by Prof. Alastair F. Donaldson and Prof. Nicolas Wu.

\section{Education}

\subsection{Imperial College London \aux{PhD in Computer Science} \rside{2022--present}}
\begin{itemize}
  \item Research on compiler testing, composable semantics, property-based testing, and the complexity of randomised testing.
  \item Full doctoral scholarship, 2022--2027.
\end{itemize}

\subsection{University of Oxford, St Anne's College \aux{MMathCompSci} \rside{2016--2020}}
\begin{itemize}
  \item Mathematics and Computer Science; first-class master's year (80.3\%).
  \item Dissertation: \textit{Systematic Pruning Methods for the Enumeration of Realistic Molecules}.
  \item Exhibition Awards in 2019 and 2020.
\end{itemize}

\section{Publications}

\begin{enumerate}
  \item \textbf{Pingshi Yu}, Chengsong Tan, Nicolas Wu, and Alastair F. Donaldson. ``Complexity Theory of Randomised Testing.'' arXiv, 2026. \href{https://arxiv.org/abs/2607.11811}{Paper \linkicon}
  \item \textbf{Pingshi Yu}, Nicolas Wu, and Alastair F. Donaldson. ``Ratte: Fuzzing for Miscompilations in Multi-Level Compilers Using Composable Semantics.'' \textit{ASPLOS}, 2025. \href{https://doi.org/10.1145/3676641.3716270}{DOI \linkicon}
  \item Mayank Sharma, \textbf{Pingshi Yu}, and Alastair F. Donaldson. ``RustSmith: Random Differential Compiler Testing for Rust.'' \textit{ISSTA}, 2023. \href{https://doi.org/10.1145/3597926.3604919}{DOI \linkicon}
  \item \textbf{Pingshi Yu}. ``Reasoning about MLIR Semantics through Effects and Handlers.'' \textit{ECOOP/ISSTA Doctoral Symposium}, 2023. \href{https://doi.org/10.1145/3597926.3605239}{DOI \linkicon}
  \item \textbf{Pingshi Yu}, Alistair J. Sterling, and Jotun Hein. ``A Novel Automated Screening Method for Combinatorially Generated Small Molecules.'' \textit{Journal of Chemical Information and Modeling}, 2021. \href{https://doi.org/10.1021/acs.jcim.0c01462}{DOI \linkicon}
\end{enumerate}

\section{Research and Industry Experience}

\subsection{Huawei Research \rside{Edinburgh, UK}}
\subsubsection{Research Intern, Programming Languages Lab \rside{2025--2026}}
\begin{itemize}
  \item Researched the implementation and application of effect handlers for Huawei's open-source Cangjie programming language.
\end{itemize}

\subsection{Credit Suisse \rside{London, UK}}
\subsubsection{Software Engineer / Quantitative Strategist \rside{2020--2022}}
\begin{itemize}
  \item Developed a foreign-exchange market-making platform and led delivery of NLP-derived market-data pipelines into its pricing engine.
  \item Built tools for quantitative research, model validation, release management, performance monitoring, and developer productivity using Python, C++, Docker, Jenkins, JavaScript, and Elasticsearch.
\end{itemize}

\subsection{Cambridge Quantum Computing (now Quantinuum) \rside{Cambridge, UK}}
\subsubsection{Research Intern \rside{2020}}
\begin{itemize}
  \item Reviewed ZX-calculus research and implemented a scalable graphical calculus in C++ for representing and simplifying large quantum circuits.
\end{itemize}

\subsection{Credit Suisse \rside{London, UK}}
\subsubsection{Technology Intern \rside{2019}}
\begin{itemize}
  \item Delivered two production projects and identified an optimisation yielding approximately $900\times$ faster bulk SQL inserts in an internal Perl API.
\end{itemize}

\subsection{Arm \rside{Cambridge, UK}}
\subsubsection{Data Science Intern \rside{2018}}
\begin{itemize}
  \item Initiated and delivered the automated data-processing backend for a company-wide capacity-planning project.
\end{itemize}

\subsection{University of Oxford, DPIR \rside{Oxford, UK}}
\subsubsection{Research Intern \rside{2017, 2018}}
\begin{itemize}
  \item Built a Python pipeline to process roughly 10,000 statutory instruments; later designed data-gathering and sentiment-analysis validation questionnaires.
\end{itemize}

\section{Selected Talks and Research Activities}

\begin{description}[style=multiline,leftmargin=4.8em]
  \item[2026] VeTSS Summer School, Exeter --- attendee.
  \item[2025] ASPLOS, Rotterdam --- presentation of \textit{Ratte}; ACE-CSR Conference, Lancaster --- attendee.
  \item[2024] ETH Z\"urich --- invited presentation; LaPRaS, London --- presentation; VeTSS Summer School, Bristol; POPL, London.
  \item[2023] ISSTA and ECOOP/ISSTA Doctoral Symposium, Seattle --- presentations; Fun in the REPL, Bristol --- presentation; Marktoberdorf Summer School; EuroLLVM Developers' Conference.
  \item[2022] Programming Languages Implementation Summer School (PLISS), Bertinoro.
\end{description}

\section{Awards and Service}

\begin{description}[style=multiline,leftmargin=5.8em]
  \item[2022--27] Full Doctoral Scholarship, Imperial College London.
  \item[2024] Poster Presentation Award, Imperial Computing Conference.
  \item[2019, 2020] Exhibition Award, St Anne's College, University of Oxford.
  \item[2026] Subreviewer, POPL.
  \item[2025] Artifact Evaluation Reviewer, FormaliSE; subreviewer, FSE.
  \item[2024] Subreviewer, ASE.
\end{description}

\section{Technical Skills and Languages}

\begin{description}[style=multiline,leftmargin=8.5em]
  \item[Programming] Haskell, Python, C++; experience with Agda, Java, Scala, Perl, and MATLAB.
  \item[Tools \& methods] Git, Unix, Docker, Jenkins, SQL, Pandas, TensorFlow, \LaTeX; compilers, property-based testing, language semantics, algorithms, statistics, and machine learning.
  \item[Languages] Chinese (native), English (native), German (basic).
  \item[Interests] Piano (Chopin and ragtime), cycling, chess, geocaching, and road trips.
\end{description}

\vspace{6pt}
\begin{center}
  {\small\color{gray}Last updated August 2026}
\end{center}

\end{document}
